\textbf{(i)} If $D(f)=D(g)$, the image of $g$ belongs to no prime of $A_f$, hence is a unit there. Likewise $f$ is a unit in $A_g$. The universal property gives canonical maps $A_f\to A_g$ and $A_g\to A_f$, whose composites are the identity because they agree with the structure map on $A$.

\textbf{(ii)} The inclusion $D(g)\subseteq D(f)$ gives $V(f)\subseteq V(g)$, so $g\in\sqrt{(f)}$ and $g^n=uf$ for some $n>0$. In $A_g$, the inverse of $f$ is $u/g^n$, so the stated formula is the unique $A$-algebra map $A_f\to A_g$. This uniqueness, together with (i), makes it independent of all choices.

\textbf{(iii)} The identity is the unique such $A$-algebra map when the opens coincide.

\textbf{(iv)} A composite of restrictions is an $A$-algebra map with the prescribed source and target, so uniqueness gives the direct restriction.

\textbf{(v)} For $\mathfrak p\in D(f)$, the canonical maps $A_f\to A_{\mathfrak p}$ are compatible and induce a map from the direct limit. Every $a/s\in A_{\mathfrak p}$ comes from $A_s$, so it is surjective. If $a/f^n$ maps to zero, choose $s\notin\mathfrak p$ with $sa=0$. Its restriction to $A_{fs}$ is then zero, and $D(fs)$ is still a neighborhood of $\mathfrak p$. Hence it represents zero in the direct limit.
